\documentclass[11pt,fontset=fandol]{ctexart}

\usepackage[a4paper,margin=1.15cm]{geometry}
\usepackage{array}
\usepackage{booktabs}
\usepackage{graphicx}
\usepackage{longtable}
\usepackage{tabularx}
\usepackage{xcolor}

\setlength{\parindent}{0pt}
\renewcommand{\arraystretch}{1.18}
\pagestyle{empty}

% 购买记录 PDF 通用模版
% 用法：
% 1. 先新建一个独立工作目录，专门放本次图片、tex 和 pdf。
% 2. 修改标题、日期。
% 3. 在首页汇总表中填入各类型条目、小计与总计。
% 4. 每个物品复制一段 \recordpage。
% 5. 有完整发票时用 \invoiceimage；缺票或未匹配时用 \invoiceplaceholder。
% 6. 最后一页 TODO 按实际缺口填写。
% 7. 已确认要入册的图片，默认先复制到这个工作目录，并改成“序号-购买记录 / 序号-消费记录 / 序号-发票”。
% 8. 这份模版默认可直接编译；示例页全部使用占位框，实际使用时再替换成真实图片路径。

\newcommand{\docTitle}{购买记录汇总}
\newcommand{\docDate}{2026年05月20日}

\newcommand{\typeRow}[4]{#1 & #2 & #3 & #4 \\}
\newcommand{\subtotalRow}[2]{\multicolumn{2}{l}{#1小计} & #2 & \\}

\newcommand{\topimage}[2]{
  \begin{minipage}[t][8.1cm][t]{0.485\linewidth}
    \centering
    \textbf{#1}\\[0.3em]
    \includegraphics[width=\linewidth,height=7.5cm,keepaspectratio]{#2}
  \end{minipage}
}

\newcommand{\topplaceholder}[2]{
  \begin{minipage}[t][8.1cm][t]{0.485\linewidth}
    \centering
    \textbf{#1}\\[0.3em]
    \fcolorbox{black}{gray!8}{
      \parbox[c][7.2cm][c]{0.92\linewidth}{
        \centering
        #2
      }
    }
  \end{minipage}
}

\newcommand{\invoiceimage}[2]{
  \begin{minipage}[t][14.6cm][t]{\linewidth}
    \centering
    \textbf{发票}\\[0.3em]
    \includegraphics[width=\linewidth,height=13.8cm,keepaspectratio,#1]{#2}
  \end{minipage}
}

\newcommand{\invoiceplaceholder}[1]{
  \fcolorbox{black}{gray!8}{
    \parbox[c][13.0cm][c]{0.96\linewidth}{
      \centering
      发票待补\\[0.6em]
      #1\\[0.4em]
      正式报销前需继续核对\\
      订单金额 = 支付金额 = 发票金额
    }
  }
}

\newcommand{\recordpage}[7]{
\clearpage
{\Large \textbf{#1}}\\[0.3em]
类型：#2 \qquad 金额：#3\\[0.7em]
\noindent
\topimage{购买记录}{#4}\hfill
\topimage{消费记录}{#5}

\vspace{0.7em}

\noindent
#6

\vfill
\textcolor{gray}{#7}
}

\newcommand{\recordpagecustom}[7]{
\clearpage
{\Large \textbf{#1}}\\[0.3em]
类型：#2 \qquad 金额：#3\\[0.7em]
\noindent
#4\hfill#5

\vspace{0.7em}

\noindent
#6

\vfill
\textcolor{gray}{#7}
}

\begin{document}

\begin{center}
  {\zihao{2}\bfseries \docTitle}\\[0.35em]
  {\large \docDate}
\end{center}

\vspace{0.25em}

{\small
\begin{center}
\begin{tabularx}{0.99\linewidth}{@{}>{\raggedright\arraybackslash}p{1.8cm} X >{\centering\arraybackslash}p{2.1cm} >{\centering\arraybackslash}p{1.8cm}@{}}
\toprule
类型 & 物品名 & 价格 & 备注 \\
\midrule
\typeRow{服装}{示例物品 A}{88.00}{}
\typeRow{服装}{示例物品 B}{128.00}{}
\subtotalRow{服装}{216.00}
\midrule
\typeRow{道具}{示例物品 C}{56.00}{}
\subtotalRow{道具}{56.00}
\midrule
\typeRow{家具}{暂无}{0.00}{}
\subtotalRow{家具}{0.00}
\midrule
\multicolumn{2}{l}{总计} & 272.00 & \\
\bottomrule
\end{tabularx}
\end{center}
}

\vspace{0.5em}

\textbf{说明：}本版按“左上购买记录、右上消费记录、下方发票”的固定构图生成。首页为分类汇总表；后续每个物品单开一页；最后一页列出缺失或未匹配材料。正式开工前，先新建一个专门工作目录，把图片、tex 和 pdf 都放进去；已确认要入册的图片，再统一改名为“序号-购买记录 / 序号-消费记录 / 序号-发票”。

\recordpagecustom
{1. 示例物品 A}
{服装}
{¥88.00}
\topplaceholder{购买记录}{把这里替换成 \texttt{\string\topimage\{购买记录\}\{./1-购买记录.png\}}}
\topplaceholder{消费记录}{把这里替换成 \texttt{\string\topimage\{消费记录\}\{./1-消费记录.png\}}}
{\invoiceplaceholder{当前未见对应完整发票。}}
{订单号：示例订单号；备注：可替换为商家或核对说明。}

\recordpagecustom
{2. 示例物品 B}
{服装}
{¥128.00}
\topplaceholder{购买记录}{把这里替换成 \texttt{\string\topimage\{购买记录\}\{./1-购买记录.png\}}}
\topplaceholder{消费记录}{把这里替换成 \texttt{\string\topimage\{消费记录\}\{./1-消费记录.png\}}}
{\invoiceplaceholder{如已有完整发票，可改为 \texttt{\string\invoiceimage\{page=1\}\{./1-发票.pdf\}}。}}
{订单号：示例订单号；发票号：如已开票再填写。}

\clearpage
{\Large \textbf{TODO}}\\[0.8em]

\textbf{缺订单截图}
\begin{itemize}
  \item 示例：某物品订单截图缺失，暂未入页。
\end{itemize}

\textbf{缺付款截图}
\begin{itemize}
  \item 示例：某物品只有订单列表页，缺付款详情页。
\end{itemize}

\textbf{缺发票}
\begin{itemize}
  \item 示例：某物品当前未见完整发票。
\end{itemize}

\textbf{已有材料但未匹配}
\begin{itemize}
  \item 示例：某发票金额与订单金额不一致，暂不入正文。
\end{itemize}

\end{document}
